\chapter{System Architecture and Design}
\label{ch:design}

\section{Overall System Architecture}
MediHaven follows a layered client--server architecture (Figure~\ref{fig:arch}). The browser renders
server-side Jinja2 templates and calls the JSON API with \texttt{fetch}; the Flask application factory
(\code{src/api/app.py}) wires security middleware and six blueprints; services hold the ML models, the
vault logic, the narrative synthesiser and the research service; SQLite, model files and the research
artifacts form the data layer.

\begin{figure}[!htbp]
\centering
\adjustbox{max width=\linewidth}{% Overall layered architecture (verified against src/api/app.py, routes/, models/, vault/, training/).
% Explicit millimetre grid: 6 columns in the upper layers, 4 in the lower layers.
\begin{tikzpicture}[x=1mm, y=1mm]
\newcommand{\band}[3]{% top y, bottom y, title
  \draw[mhgray!60, dashed, rounded corners=3pt] (0,#1) rectangle (170,#2);
  \node[layer label] at (1.5,#1-0.8) {#3};}
\band{0}{-19}{Presentation: Jinja2 templates, JavaScript, Chart.js (browser)}
\foreach \x/\t in {15/Landing \& Login, 43/Physician Triage, 71/Patient Vault, 99/QR Scanner,
                   127/Admin Overview, 155/Research (opt-in)}
  \node[blue box, minimum width=25mm, text width=23mm] at (\x,-11.5) {\t};

\band{-25}{-58}{Application: Flask app factory served by Waitress (wsgi.py) or run\_demo.py}
\node[teal box, minimum width=164mm] at (85,-36)
  {session RBAC $\cdot$ CSP and security headers $\cdot$ rate limits $\cdot$ research-mode boundary $\cdot$ JSON error handlers};
\foreach \x/\t in {15/auth, 43/patients, 71/predictions, 99/alerts, 127/vault, 155/research}
  \node[teal box, minimum width=25mm] at (\x,-49.5) {/api/\t};

\band{-64}{-86}{Services}
\node[amber box, minimum width=38mm, text width=35mm] (ens) at (22,-76.5) {4-model ensemble\\K-Means $\cdot$ DT $\cdot$ KNN $\cdot$ MLP};
\node[amber box, minimum width=38mm, text width=35mm] (qr) at (64,-76.5) {Vault \& QR service\\HMAC-signed JWT passes};
\node[amber box, minimum width=38mm, text width=35mm] (nlp) at (106,-76.5) {Narrative synthesiser\\de-identified prompts};
\node[amber box, minimum width=38mm, text width=35mm] (rs) at (148,-76.5) {Research service\\verified artifacts};

\band{-92}{-114}{Data and external services}
\node[store, minimum width=38mm, text width=35mm] (ms) at (22,-104.5) {models\_store/\\joblib model files};
\node[store, minimum width=38mm, text width=35mm] (db) at (64,-104.5) {SQLite (WAL mode)\\8 tables, 8 indexes};
\node[store, minimum width=38mm, text width=35mm] (gm) at (106,-104.5) {Google Gemini API\\external, optional};
\node[store, minimum width=38mm, text width=35mm] (ra) at (148,-104.5) {Research artifacts\\outside repo \& OneDrive};

\draw[flow] (85,-19) -- node[lbl, right=1mm]{HTTP + JSON (fetch), session cookie} (85,-25);
\draw[flow] (85,-58) -- (85,-64);
\draw[flow] (ens.south) -- (ms.north);
\draw[flow] (ens.south east) -- (db.north west);
\draw[flow] (qr.south) -- (db.north);
\draw[dflow] (nlp.south) -- node[lbl, right]{synthetic mode only} (gm.north);
\draw[flow] (rs.south) -- node[lbl, right]{hash-verified} (ra.north);
\end{tikzpicture}
}
\caption{Layered architecture of MediHaven}
\label{fig:arch}
\end{figure}

\section{Component and Module Description}
\begin{table}[!htbp]
\centering
\caption{Major modules: purpose, inputs, outputs and dependencies}
\label{tab:modules}
\small
\begin{tabularx}{\textwidth}{@{}p{40mm}YYp{30mm}@{}}
\toprule
Module & Purpose & Input $\rightarrow$ output & Depends on\\
\midrule
\code{api/app.py}, \texttt{auth.py} & app factory, RBAC, headers, rate limits & HTTP request $\rightarrow$ routed response & config, security\\
\texttt{routes/patients, predictions, alerts} & triage and charts & patient id $\rightarrow$ JSON packet & ensemble, SQLite\\
\code{models/ensemble.py} + 4 models & synthetic risk inference & 30 features $\rightarrow$ tier, score, votes & scikit-learn\\
\code{vault/*}, \code{routes/vault} & records, QR passes, audit & scope, token $\rightarrow$ scoped records & PyJWT, OpenCV\\
\code{nlp/summarizer.py} & narratives & coarse model outputs $\rightarrow$ text & Gemini API (optional)\\
\code{tasks/build_mimic_t1.py} & research dataset & MIMIC CSV $\rightarrow$ hashed parquet & DuckDB, pandas\\
\code{training/*} & research training, verification, serving & dataset $\rightarrow$ verified run, OOF view & scikit-learn\\
\texttt{bootstrap.py}, \texttt{wsgi.py} & first start, production entry & config $\rightarrow$ ready app & Waitress\\
\bottomrule
\end{tabularx}
\end{table}

\section{Architecture Selection and Justification}
A layered, modular monolith was chosen over microservices: the workload is small (Chapter~\ref{ch:scal}),
the team is three students, models load once into process memory (no network hop per prediction), and a
single SQLite file keeps deployment trivial. Blueprints and service classes keep the layers separable, so
the data layer or the model service could later be extracted.

\section{Data Flow Diagrams}
Figure~\ref{fig:dfd0} shows the context and Figure~\ref{fig:dfd1} the level-1 processes and stores. The
Gemini flow is dashed because it is optional and only ever receives coarse model outputs.

\begin{figure}[!htbp]
\centering
\adjustbox{max width=0.85\linewidth}{% Context-level (level 0) DFD: one process, external entities, labelled flows.
\begin{tikzpicture}[x=1mm, y=1mm]
\node[process, minimum size=26mm, font=\footnotesize] (sys) at (85,0) {\textbf{0}\\MediHaven\\system};
\node[box, minimum width=28mm] (phys) at (14,24) {Physician};
\node[box, minimum width=28mm] (pat) at (14,-24) {Patient};
\node[box, minimum width=28mm] (adm) at (156,24) {Administrator};
\node[box, minimum width=28mm] (gem) at (156,-24) {Google Gemini API\\(synthetic mode only)};
\node[box, minimum width=40mm] (mimic) at (85,-38) {MIMIC-IV Demo files\\(offline research build)};
% each actor: request flow above, response flow below
\draw[flow] ($(phys.east)+(0,1.5)$) -- node[lbl, above, sloped, pos=0.45]{credentials, chart requests, QR scans} (sys.130);
\draw[flow] (sys.155) -- node[lbl, below, sloped, pos=0.55]{tiers, charts, vault records} ($(phys.east)+(0,-1.5)$);
\draw[flow] ($(pat.east)+(0,-1.5)$) -- node[lbl, below, sloped, pos=0.45]{records, share scope} (sys.230);
\draw[flow] (sys.205) -- node[lbl, above, sloped, pos=0.55]{QR pass, access log} ($(pat.east)+(0,1.5)$);
\draw[flow] ($(adm.west)+(0,1.5)$) -- node[lbl, above, sloped, pos=0.45]{reassign, dispatch} (sys.50);
\draw[flow] (sys.25) -- node[lbl, below, sloped, pos=0.55]{census, model telemetry} ($(adm.west)+(0,-1.5)$);
\draw[dflow] (sys.335) -- node[lbl, above, sloped, pos=0.55]{de-identified prompt} ($(gem.west)+(0,1.5)$);
\draw[dflow] ($(gem.west)+(0,-1.5)$) -- node[lbl, below, sloped, pos=0.35]{narrative text} (sys.310);
\draw[flow] (mimic.north) -- node[lbl, left]{stay-level features} (sys.south);
\end{tikzpicture}
}
\caption{Context-level (level-0) data flow diagram}
\label{fig:dfd0}
\end{figure}

\begin{figure}[!htbp]
\centering
\adjustbox{max width=0.88\linewidth}{% Level-1 DFD: processes 1.0-5.0, data stores D1-D3 and external entities.
\begin{tikzpicture}[x=1mm, y=1mm]
\node[box, minimum width=24mm] (user) at (10,0) {Physician /\\Patient /\\Admin};
\node[process] (p1) at (48,24) {\textbf{1.0}\\Authenticate\\\& authorise};
\node[process] (p2) at (48,-24) {\textbf{2.0}\\Triage \&\\predict};
\node[process] (p3) at (98,24) {\textbf{3.0}\\Vault \&\\QR sharing};
\node[store, minimum width=30mm] (d1) at (98,0) {D1\;SQLite database};
\node[process] (p4) at (98,-24) {\textbf{4.0}\\Narrative\\synthesis};
\node[store, minimum width=30mm] (d2) at (10,-42) {D2\;models\_store};
\node[box, minimum width=24mm] (ext) at (98,-48) {Gemini API};
\node[process] (p5) at (150,-12) {\textbf{5.0}\\Research\\pipeline};
\node[store, minimum width=32mm] (d3) at (150,18) {D3\;research artifacts};
\node[box, minimum width=28mm] (mim) at (150,-42) {MIMIC-IV Demo};
\draw[flow] (user) -- node[lbl, above, sloped]{credentials} (p1);
\draw[flow] (user) -- node[lbl, below, sloped]{chart request} (p2);
\draw[flow] (p1) -- node[lbl, left]{session role} (p2);
\draw[flow] (p1) -- node[lbl, above]{session role} (p3);
\draw[flow] (p3) -- node[lbl, right]{tokens, access log} (d1);
\draw[flow] (d1) -- node[lbl, above, sloped]{vitals, labs} (p2);
\draw[flow] (d2) -- node[lbl, below, sloped]{fitted models} (p2);
\draw[flow] (p2) -- node[lbl, below]{risk packet} (p4);
\draw[dflow] (p4) -- node[lbl, right]{coarse outputs only} (ext);
\draw[flow] (mim) -- node[lbl, right]{CSV} (p5);
\draw[flow] (p5) -- node[lbl, right]{verified run} (d3);
\end{tikzpicture}
}
\caption{Level-1 data flow diagram}
\label{fig:dfd1}
\end{figure}

\section{UML Diagrams}
The use-case diagram appears in Figure~\ref{fig:usecase}. Figure~\ref{fig:class} shows the main classes:
the ensemble aggregates four model wrappers, and the research service composes the verified artifacts and
the lazily built stay view. Figure~\ref{fig:seq} traces the redemption of a QR pass; the use count is
claimed with a single conditional \texttt{UPDATE}, so concurrent scans cannot overspend a pass.
Figure~\ref{fig:activity} gives the activity flow of a prediction, Figure~\ref{fig:state} the life cycle
of a pass, and Figure~\ref{fig:deploy} the deployment options.

\begin{figure}[!htbp]
\centering
\adjustbox{max width=\linewidth}{% Class diagram of the main model and service classes (names and members from the source code).
\begin{tikzpicture}[x=1mm, y=1mm,
  cls/.style={draw=mhgray, fill=white, rectangle split, rectangle split parts=3, align=left,
              font=\scriptsize, inner sep=2pt, text width=#1, rectangle split part fill={mhblue!12,white,white}},
  cls/.default=40mm]
\node[cls=58mm] (ens) at (85,32) {\textbf{EnsembleClinicalPredictor}
  \nodepart{second} weights: dict (DT .30, MLP .30, KNN .25, K-Means .15)\\ scaler: StandardScaler
  \nodepart{third} train\_from\_data() / load()\\ predict\_patient(id): dict};
\node[cls=39mm] (km) at (18,4) {\textbf{KMeansRiskClustering}
  \nodepart{second} n\_clusters = 4\\ cluster\_to\_tier\_code
  \nodepart{third} fit(), predict\_proba()};
\node[cls=39mm] (dt) at (62,4) {\textbf{DecisionTreeRiskClassifier}
  \nodepart{second} max\_depth = 5
  \nodepart{third} fit(), predict(), explain()};
\node[cls=39mm] (knn) at (106,4) {\textbf{KNNSimilarityMatcher}
  \nodepart{second} n\_neighbors = 5
  \nodepart{third} find\_similar(), recommend\_protocol()};
\node[cls=39mm] (nn) at (150,4) {\textbf{NeuralNetworkRiskModel}
  \nodepart{second} hidden\_layer\_sizes = (64, 32)
  \nodepart{third} fit(), predict\_proba()};
\foreach \c in {km,dt,knn,nn} \draw[mhgray, -{Diamond[open,length=3mm]}] (\c.north) -- (ens.south);

\node[cls=48mm] (rms) at (30,-30) {\textbf{ResearchModelService}
  \nodepart{second} artifacts: VerifiedArtifacts\\ stays: LazyStayView
  \nodepart{third} status(), dashboard(), predict()\\ stay\_roster(), stay\_detail(n)};
\node[cls=40mm] (va) at (88,-30) {\textbf{VerifiedArtifacts}
  \nodepart{second} run\_manifest, metrics\\ features (74), models
  \nodepart{third} verify\_and\_load(dir, pin)};
\node[cls=38mm] (lsv) at (145,-30) {\textbf{LazyStayView}
  \nodepart{second} lock, cached view or error
  \nodepart{third} get(): StayView};
\node[cls=44mm] (sdt) at (88,-58) {\textbf{SmoothedDecisionTree}
  \nodepart{second} max\_depth 3, min leaf 8, $m = 5$\\ leaf\_rate\_ (m-estimate)
  \nodepart{third} fit(), predict\_proba()};
\node[cls=48mm] (sum) at (30,-58) {\textbf{ClinicalNarrativeSummarizer}
  \nodepart{second} GEMINI\_MODELS, api\_key
  \nodepart{third} generate(), external\_calls\_allowed()};
\draw[mhgray, -{Diamond[length=3mm]}] (va.west) -- (rms.east);
\draw[mhgray, -{Diamond[length=3mm]}] (lsv.north west) to[bend right=12] (rms.north east);
\draw[dflow] (va.south) -- node[lbl, right]{loads pipeline using} (sdt.north);
\end{tikzpicture}
}
\caption{Class diagram of the model and service layer}
\label{fig:class}
\end{figure}

\begin{figure}[!htbp]
\centering
\adjustbox{max width=\linewidth}{% Sequence diagram: clinician redeems a QR access pass (src/api/routes/vault.py, src/vault/*).
\begin{tikzpicture}[x=1mm, y=1mm]
\foreach \x/\n/\t in {8/c/Clinician browser, 44/r/vault route (POST), 82/a/access\_controller,
                      118/s/qr\_scanner, 152/d/SQLite}
  {\node[box, minimum width=28mm, text width=26mm] (\n) at (\x,0) {\t};
   \draw[mhgray, dashed] (\x,-4) -- (\x,-80);}
\newcommand{\msg}[5][]{\draw[flow,#1] (#2,#4) -- node[lbl, above]{#5} (#3,#4);}
\msg{8}{44}{-10}{token or image (session cookie)}
\node[lbl, right, align=left, font=\scriptsize\itshape] at (45,-17) {RBAC, payload type check,\\signed-in user = auditor};
\msg{44}{82}{-24}{access\_vault\_with\_token()}
\msg{82}{118}{-30}{scan\_and\_validate\_qr()}
\msg{118}{152}{-36}{token by hash: revoked? expiry? uses?}
\msg[dflow]{152}{118}{-42}{token record}
\node[lbl, align=center] at (135,-48) {HMAC-SHA256 JWT check\\(signature, expiry)};
\msg[dflow]{118}{82}{-54}{valid / reason}
\msg{82}{152}{-60}{atomic UPDATE use\_count (if uses remain)}
\msg{82}{152}{-66}{INSERT vault\_access\_log}
\msg[dflow]{82}{44}{-72}{scoped records}
\msg[dflow]{44}{8}{-78}{200 JSON or 403 reason}
\end{tikzpicture}
}
\caption{Sequence diagram: redeeming a QR access pass}
\label{fig:seq}
\end{figure}

\begin{figure}[!htbp]
\centering
\adjustbox{max width=0.75\linewidth}{% Activity diagram: synthetic risk prediction (predictions.py -> ensemble.predict_patient).
\begin{tikzpicture}[node distance=5mm, act/.style={box, text width=100mm, font=\footnotesize}]
\node[term] (s) {start: GET /api/patients/\{id\}/predict};
\node[decision, below=of s] (auth) {session role\\physician or admin?};
\node[act, below=of auth] (load) {load vitals, labs and medications from SQLite; engineer 30 features (latest values,\\72-hour means and OLS slopes, shock index, MAP)};
\node[act, below=of load] (scale) {scale the feature vector with the stored StandardScaler};
\node[act, below=of scale] (four) {K-Means, decision tree, KNN and MLP each return a 4-tier probability vector};
\node[act, below=of four] (fuse) {fuse: $p = 0.30\,p_{DT} + 0.30\,p_{MLP} + 0.25\,p_{KNN} + 0.15\,p_{KM}$; tier $= \arg\max_t p_t$;\\score $= \sum_t p_t s_t$ with $s = (0.10, 0.45, 0.78, 0.94)$};
\node[act, below=of fuse] (save) {store a prediction row; return tier, score, per-model votes, rule path and similar cases};
\node[term, below=of save] (e) {end: 200 JSON};
\node[term, fill=mhred!12, right=14mm of auth] (deny) {401 / 403};
\node[term, fill=mhred!12, right=6mm of scale.east, anchor=west, yshift=-6mm] (fail) {500, generic message};
\draw[flow] (s) -- (auth);
\draw[flow] (auth) -- node[lbl, right]{yes} (load);
\draw[flow] (auth) -- node[lbl, above]{no} (deny);
\draw[flow] (load) -- (scale); \draw[flow] (scale) -- (four); \draw[flow] (four) -- (fuse);
\draw[flow] (fuse) -- (save); \draw[flow] (save) -- (e);
\draw[dflow] (four.east) -| node[lbl, below, pos=0.25]{model error} (fail.south);
\end{tikzpicture}
}
\caption{Activity diagram: synthetic risk prediction}
\label{fig:activity}
\end{figure}

\begin{figure}[!htbp]
\centering
\adjustbox{max width=0.85\linewidth}{% State machine of a QR access pass (qr_generator.py, qr_scanner.py, access_controller.py).
\begin{tikzpicture}[x=1mm, y=1mm]
\fill[mhgray] (0,0) circle (1.6);
\node[state] (act) at (34,0) {Active\\\scriptsize use\_count $<$ max\_uses};
\node[state] (used) at (100,14) {Exhausted\\\scriptsize use\_count = max\_uses};
\node[state] (exp) at (100,0) {Expired\\\scriptsize now $>$ expires\_at};
\node[state] (rev) at (100,-14) {Revoked\\\scriptsize by the patient};
\draw[flow] (1.6,0) -- node[lbl, above]{share (patient)} (act);
\draw[flow] (act) to[loop above] node[lbl, above]{granted scan: atomic use\_count + 1} (act);
\draw[flow] (act) -- node[lbl, above, sloped]{last allowed scan} (used);
\draw[flow] (act) -- node[lbl, above]{TTL elapses ($\le$ 30 days)} (exp);
\draw[flow] (act) -- node[lbl, below, sloped]{revoke} (rev);
\node[font=\scriptsize, align=left, anchor=west] at (122,0) {terminal: every further\\scan is denied and\\written to the access log};
\end{tikzpicture}
}
\caption{State machine of a QR access pass}
\label{fig:state}
\end{figure}

\begin{figure}[!htbp]
\centering
\adjustbox{max width=\linewidth}{% Deployment diagram: public synthetic deployment vs local-only research mode (wsgi.py, Dockerfile, run_demo.py).
\begin{tikzpicture}[x=1mm, y=1mm]
\node[box, minimum width=24mm] (user) at (12,12) {Browser\\(HTTPS)};
\node[box, minimum width=24mm, text width=22mm] (proxy) at (12,-12) {TLS reverse\\proxy};
\draw[mhblue, rounded corners=4pt, thick] (28,30) rectangle (104,-30);
\node[font=\scriptsize\bfseries, text=mhblue, anchor=north west, text width=72mm] at (29.5,29) {Server or container (synthetic demo, production mode)};
\node[teal box, text width=30mm] (w) at (48,12) {Waitress, 8 threads\\wsgi.py $\rightarrow$ Flask};
\node[amber box, text width=26mm] (b) at (87,12) {src.bootstrap:\\seed, train if missing};
\node[store, text width=30mm] (db) at (48,-14) {SQLite volume\\/app/database};
\node[store, text width=26mm] (ms) at (87,-14) {models\_store,\\data/processed};
\draw[flow] (user) -- (proxy); \draw[flow] (proxy.east) -- (w.south west);
\draw[flow] (w) -- (db); \draw[flow] (b) -- (ms); \draw[flow] (b.west) -- (w.east);
\draw[mhred, rounded corners=4pt, thick] (110,30) rectangle (170,-30);
\node[font=\scriptsize\bfseries, text=mhred, anchor=north west, text width=56mm] at (111.5,29) {Researcher's machine only (research mode)};
\node[teal box, text width=50mm] (rd) at (140,12) {run\_demo.py on 127.0.0.1\\(loopback and strong secrets enforced)};
\node[store, text width=50mm] (ra) at (140,-12) {C:\textbackslash research: MIMIC-IV Demo, dataset,\\verified run (outside Git and OneDrive)};
\draw[flow] (rd) -- (ra);
\node[font=\scriptsize, text=mhred] at (140,-26) {wsgi.py refuses research mode};
\end{tikzpicture}
}
\caption{Deployment diagram: production (synthetic) and local research mode}
\label{fig:deploy}
\end{figure}

\section{Database Design}
The schema (\code{database/schema.sql}, Figure~\ref{fig:er}) has eight tables in two domains: clinical
(patients, vitals, lab results, medications, predictions, outcomes) and vault (records, access tokens,
access log). Every table has a surrogate integer primary key; every child table references
\texttt{patients(patient\_id)} with \texttt{ON DELETE CASCADE}, and the access log also references the token.
Each non-key attribute depends only on its table's key, so the schema is in third normal form; JSON
columns (\code{scope_json}, \code{structured_data}) hold small, opaque lists that are never queried
by element.

\begin{figure}[!htbp]
\centering
\adjustbox{max width=\linewidth}{% ER diagram of database/schema.sql (8 tables; PK underlined, FK in italics; all FKs ON DELETE CASCADE).
\begin{tikzpicture}[x=1mm, y=1mm,
  tbl/.style={entity, text width=#1}, tbl/.default=38mm]
\node[tbl=44mm] (pat) at (85,0) {\textbf{patients}
  \nodepart{second}\underline{patient\_id}\\ mrn (UNIQUE)\\ full\_name, age (0--130), gender\\ admission\_date, status, ward\\ bed\_number, assigned\_physician};
\node[tbl] (vit) at (18,24) {\textbf{vitals}
  \nodepart{second}\underline{vital\_id}, \textit{patient\_id}\\ recorded\_at, heart\_rate\\ BP sys/dia, temperature\\ SpO$_2$, resp.\ rate (CHECKs)};
\node[tbl] (lab) at (18,-8) {\textbf{lab\_results}
  \nodepart{second}\underline{lab\_id}, \textit{patient\_id}\\ recorded\_at, glucose\\ creatinine, wbc\_count\\ bun, hemoglobin, platelets};
\node[tbl] (med) at (18,-36) {\textbf{medications}
  \nodepart{second}\underline{medication\_id}, \textit{patient\_id}\\ drug\_name, dosage, route (CHECK)\\ frequency, start/end, is\_active};
\node[tbl] (pred) at (152,24) {\textbf{predictions}
  \nodepart{second}\underline{prediction\_id}, \textit{patient\_id}\\ predicted\_at, risk\_tier\\ risk\_score [0,1], model\_used\\ rule path, protocol};
\node[tbl] (out) at (152,-6) {\textbf{outcomes}
  \nodepart{second}\underline{outcome\_id}, \textit{patient\_id}\\ primary\_diagnosis\\ readmitted\_30d, icu\_transfer};
\node[tbl] (vlt) at (85,-36) {\textbf{medical\_vault}
  \nodepart{second}\underline{vault\_id}, \textit{patient\_id}\\ category (CHECK), title\\ structured\_data, is\_sensitive};
\node[tbl=40mm] (tok) at (152,-36) {\textbf{vault\_access\_tokens}
  \nodepart{second}\underline{token\_id}, \textit{patient\_id}\\ token\_hash (UNIQUE)\\ scope\_json, expires\_at\\ max\_uses, use\_count, revoked};
\node[tbl=40mm] (log) at (152,-66) {\textbf{vault\_access\_log}
  \nodepart{second}\underline{log\_id}, \textit{token\_id}, \textit{patient\_id}\\ accessed\_by, access\_status\\ accessed\_at, accessed\_scope};
\newcommand{\rel}[3]{\draw[mhgray, thick] (#1) -- node[lbl, pos=0.12]{1} node[lbl, pos=0.86]{N} (#2);}
\rel{pat.north west}{vit.east}{}
\rel{pat.west}{lab.east}{}
\rel{pat.south west}{med.east}{}
\rel{pat.north east}{pred.west}{}
\rel{pat.east}{out.west}{}
\rel{pat.south}{vlt.north}{}
\rel{pat.south east}{tok.north west}{}
\rel{tok.south}{log.north}{}
\draw[mhgray, thick] (pat.south east) ++(0,-2) to[out=-40, in=180] node[lbl, pos=0.85]{N} (log.west);
\end{tikzpicture}
}
\caption{Entity--relationship diagram (PK underlined, FK in italics; 1:N relationships)}
\label{fig:er}
\end{figure}

\textbf{Integrity constraints.} \texttt{CHECK} constraints bound physiology (e.g.\ heart rate 20--260 bpm,
SpO$_2$ 50--100\,\%), enumerations (gender, status, route, vault category, access status) and probabilities
($0 \le$ risk score $\le 1$); \texttt{mrn} and \code{token_hash} are \texttt{UNIQUE}.
\textbf{Indexing.} Eight composite indexes follow the access paths: time series by
\texttt{(patient\_id, recorded\_at DESC)} for vitals and labs, census by \texttt{(status, ward)}, vault by
\texttt{(patient\_id, category)}, token lookup by \texttt{(token\_hash, revoked, expires\_at)}, and the
access log by patient and by token with \texttt{accessed\_at DESC}. \textbf{Representative data.} On
26~October 2026 the local demo database held 100 patients, 4{,}368 vitals rows, 200 lab rows, 255
medications, 338 vault records and 210 issued passes; for example patient~1 (MRN-2026-0001, 53 years, M,
Surgical Recovery, Bed-301) had a latest reading of HR 65.1 bpm, BP 121/80 mmHg, SpO$_2$ 98.5\,\%. All of
these are synthetic.

\section{Security and Error Handling Design}
\label{sec:secdesign}
\begin{itemize}
\item \textbf{Authentication:} only password hashes (scrypt, Werkzeug) are stored; in production, admin
      accounts need a hash supplied through the environment and the public demo logins need
      \texttt{DEMO\_LOGINS=1}; login is limited to 10 attempts per 5 minutes per client.
\item \textbf{Sessions:} signed cookies, HttpOnly, SameSite=Lax, Secure in production, 8-hour lifetime;
      the application refuses default, placeholder or short secrets in production, on any non-loopback
      host and in research mode.
\item \textbf{Authorisation:} role checks on every page and API; patients can reach only their own vault;
      audit entries take the clinician's identity from the session, never from the request body.
\item \textbf{Input validation:} JSON object bodies only; bounded integers (pass lifetime $\le$ 30 days,
      uses $\le$ 1000); scanner payloads must be a token or an image data URI no larger than 4096$\times$4096
      pixels; requests above 8~MB are rejected (413).
\item \textbf{Browser hardening:} Content-Security-Policy allowing only same-origin scripts (Chart.js is
      self-hosted, no inline scripts), \texttt{X-Frame-Options: DENY}, \texttt{nosniff}, HSTS over HTTPS,
      \texttt{no-store} on API responses.
\item \textbf{Errors and logging:} JSON error envelopes with generic messages; details are logged
      server-side only; logging falls back to the console when the log folder is not writable.
\item \textbf{Sensitive data:} Gemini receives only coarse outputs with a \texttt{[PATIENT]}
      placeholder over verified TLS; research mode never calls it, refuses non-loopback binding, verifies
      SHA-256 hashes before deserialising any model, and never returns MIMIC identifiers or dates.
\end{itemize}
